\section{Reproducibility and Artifact}
\label{sec:reproducibility-artifact}

\subsection{Artifact scope}
\label{subsec:artifact-scope}

The reproducibility artifact preserves the evidence path from the executable
protocol and formal models to the derived results reported in this paper,
rather than only source code or final tables.

The scientific-source layer contains the Java cross-shard implementation,
the TLA+ specifications and configurations used for bounded model checking,
and the Alloy models and commands used for bounded relational analysis.
These are the executable and formal objects evaluated by the campaign.

The frozen experimental-definition layer records the protocol version,
research questions and hypotheses, configuration families, bound profiles,
mutation cases, trace scenarios, deterministic seeds, repetitions, resource
limits, tool configurations, and rules for incomplete executions, separately
from the observed outcomes.

The evidence layer retains raw task results, derived datasets, and generated
tables and figures, preserving the link from reported summaries to the
underlying definitive observations.

The reproduction layer provides scripts and targets for structural checks,
isolated workspace preparation, representative verification and
trace-conformance smoke cases, analysis regeneration from preserved raw
results, and comparison with the reference artifact.

Integrity metadata, including manifests, content hashes, source-version
identifiers, and reproduction reports, connects the preserved campaign with
regenerated outputs and makes unexpected changes observable.

Together, the artifact supports inspection, representative partial
re-execution, analysis regeneration, and integrity verification. The
independent workflow does not require re-executing every definitive task and
distinguishes re-execution from regeneration based on preserved raw evidence.

\subsection{Reproduction workflow}
\label{subsec:reproduction-workflow}

The workflow separates structural validation, selective re-execution,
analysis regeneration, and integrity comparison.

First, it validates that the scientific sources, frozen experimental
definition, preserved raw results, derived outputs, and reproduction scripts
are available.

Second, it prepares an isolated workspace, records the source revision, and
creates regenerated-output directories without modifying the reference
campaign.

The third stage executes a representative smoke suite. The suite covers
a valid and a mutated TLC configuration, a valid and a mutated Alloy
configuration, one valid implementation trace, and one deliberately
corrupted trace. These executions exercise the principal formal and
implementation-facing paths of the artifact without repeating the full
definitive campaign.

The fourth stage regenerates the paper-level analysis from the preserved
raw experimental results. Derived datasets, tables, figures, and
research-question summaries are reconstructed using the same analysis
pipeline used for the reference artifact.

Fifth, regenerated outputs are compared with preserved references using
content hashes and reproduction gates to detect unexpected differences.

A reproduction verdict is issued only after structural checks, representative
executions, analysis regeneration, and integrity comparisons have completed.

Representative re-execution checks that the principal executable paths remain
operational, whereas analysis regeneration checks that paper-level evidence
can be reconstructed from the preserved definitive observations.

Neither operation is a new independent measurement of all definitive tasks.
Preserved raw results remain the source observations, while the smoke suite
only exercises selected artifact paths.

\subsection{Independent reproduction evidence}
\label{subsec:independent-reproduction}

The artifact was evaluated through an independent reproduction attempt using
a separate operating-system user, clone, and execution workspace from those
used to construct the reference campaign.

The reproduction used source revision
\texttt{6cd88c37\allowbreak 7afd23fe\allowbreak e4998882\allowbreak f91142d7\allowbreak 1e7d963e}
and packaged-artifact SHA-256 digest
\texttt{d464888e\allowbreak 9f3e5d8c\allowbreak c64ef5d2\allowbreak 2cc7b7c2\allowbreak 4f83e385\allowbreak 3f5825f1\allowbreak 8f23de26\allowbreak adf6a6e6}.

The reproduction workflow completed all ten predefined validation gates.
These gates covered artifact structure, workspace preparation,
representative formal and trace-conformance executions, analysis
regeneration, and integrity comparison. No reproduction incident was
recorded during the successful attempt.

All 32 expected output hashes matched their preserved references. The
successful reproduction therefore obtained 10/10 completed gates and 32/32
matching analysis artifacts, with zero recorded incidents.

The executable part of the reproduction consisted of the representative
smoke suite described above rather than a re-execution of the complete
definitive campaign. The full set of paper-level derived outputs was
instead regenerated from the preserved definitive raw observations and
checked against the reference artifact.

The attempt therefore demonstrates representative executable-path validation
and complete analysis regeneration from preserved raw evidence, not an
independent repetition of every definitive experimental measurement.

The reproduction provides process separation but not hardware independence:
a separate operating-system user, clone, and workspace were used on the same
native Linux host. This reduces dependence on the original working tree and
user-local execution state, while a second physical machine remains an
additional possible validation step.

\subsection{Integrity and traceability}
\label{subsec:artifact-integrity}

Source-version identifiers link the implementation, formal models, and
reproduction scripts to a repository state, while the frozen experimental
definition identifies the campaign associated with the reported results.

Preserved raw results are the primary empirical record. Derived datasets,
research-question summaries, tables, and figures are generated from these
observations through the documented analysis pipeline.

Integrity manifests record cryptographic hashes for analysis artifacts
expected to remain unchanged. Regenerated files are hashed again and compared
with the preserved references.

A matching hash establishes content identity for the corresponding
artifact. It does not independently establish scientific correctness,
but it provides a deterministic check that the reproduction pipeline
reconstructed the expected analysis output from the preserved evidence.

A hash mismatch is an observable reproduction difference requiring
investigation. The reproduction report records gate outcomes, integrity
comparisons, and incidents.

The resulting traceability chain can be summarized as:

\begin{quote}
source revision $\rightarrow$
frozen experimental protocol $\rightarrow$
preserved raw observations $\rightarrow$
derived analysis artifacts $\rightarrow$
integrity manifest $\rightarrow$
reproduction comparison.
\end{quote}

This chain supports bidirectional inspection, from paper-level values back to
preserved observations and from regenerated artifacts to the reference
manifest.

\subsection{Reproducibility boundaries}
\label{subsec:reproducibility-boundaries}

The successful reproduction was not a full experimental replication. It did
not re-execute every definitive task, and preserved raw observations remained
the empirical source for complete analysis regeneration.

It also does not establish hardware-independent reproducibility. The separate
user, clone, and workspace remained on the same native Linux host. A second
physical machine would provide a stronger test of environment portability.

Similarly, the artifact does not claim universal reproducibility across
arbitrary operating systems, processors, formal-tool versions, Java
runtimes, or resource configurations. Changes in these components may
affect execution behavior, verification cost, or generated outputs.

Finally, matching hashes establish content correspondence, not scientific
validity. The conclusions remain dependent on the experimental design,
formal abstractions, statistical treatment, and stated claim boundaries.
